\documentclass[a4paper,11pt]{article}
\usepackage{../themeUGM/soaltugas}
\usepackage[T1]{fontenc}
\usepackage{tgbonum}
\usepackage{hyperref}
% Metadata
\date{\today}
\setmodule{Model Relasional dan Querynya}
\setterm{Semester 1, 2026}

%-------------------------------%
% Other details
% TODO: Fill these
%-------------------------------%
\title{Tugas Kelompok: I}
\setfacultyname{Fakultas MIPA}
\setdeptname{S1 Ilmu Komputer}
\setcoursename{Basis Data}

%-------------------------------%
% Add / Delete commands and packages
% TODO: Add / Delete here as you need
%-------------------------------%
\usepackage{amsmath,amssymb,bm}
\newcommand{\tekskode}[1]{{\bf\large\texttt{#1}}}
% Main document
\begin{document}
    % Add header
    \header{}
\fontfamily{qpl} \selectfont
\begin{deliverablebox}{Tugas terdiri atas empat bagian}
  \begin{enumerate}
    \itemsep0em
    \item menerjemahkan batasan dan proses bisnis ke dalam diagram ER (\textit{Entity Relationship Diagram}). Usulan disertai penjelasan proses pembentukan ERD
     \item merancang dan mengusulkan skema basis data relasional yang sesuai dengan ERD yang diusulkan
     \item Jika diberikan query, mengusulkan  representasi query dalam  ekspresi aljabar relasional
     \item Mengusulkan penyelesaian pertanyaan dalam bentuk query SQL. Bila ada lebih dari satu kemungkinan query (dengan operasi atau klausa yang berbeda), query alternatif juga perlu dituliskan.
\end{enumerate}  
\end{deliverablebox}

\begin{infobox}{Ketentuan Tugas}
\begin{enumerate}
    \itemsep0em
     \item Pekerjaan dapat dikerjakan secara berkelompok dengan anggota kelompok \textbf{maksimal 4 orang}.
     \item Setiap anggota kelompok \textbf{harus} berkontribusi secara memadai, dan menguasai materi yang dihasilkan oleh kelompoknya.
     \item Pekerjaan harus disajikan secara \textit{jelas, lengkap dan rapi, termasuk penggunaan bahasa formal yang mudah dipahami}
     \item Pekerjaan harus dari pemikiran anggota tim, dan disepakati oleh semua anggota
\end{enumerate}  
\end{infobox}
\begin{warnbox}{Petunjuk Pengerjaan}
\begin{enumerate}
    \itemsep0em
    \item Tugas dikerjakan dengan menggunakan pengolah teks \LaTeX dengan template yang diberikan terlampir pada penugasan ini
    \item Cara penggunaan Template \LaTeX:
    \begin{enumerate}
    \itemsep0em
     \item Template dapat dijalankan pada platform \textbf{offline} ataupun \textbf{online}.
     \item Untuk platform offline, Anda dapat menggunakan \textit{tools}  Windows dengan \textbf{MikTeX}, Linux dengan \textbf{TeXLive} atau Mac dengan \textbf{MacTeX}
     \item Untuk platfom online, Anda dapat menggunakan layanan 
     \begin{itemize}
         \item \href{https://www.overleaf.com/}{Overleaf}
         \item \href{https://www.texpage.com/}{TexPage}
     \end{itemize}
     Caranya dengan membuat project baru, lalu memilih fitur \textbf{Existing Project} (pada Overleaf), atau \textbf{Upload Project} (pada Texpage)
     \item Anda yang kelas IUP dapat melihat detail tugas dengan menjalankan/mengkompilasi file \tekskode{AssignmentTemplate.tex} yang ada pada folder \tekskode{IUP}. 
     
     Anda yang dari non IUP dapat mengkompilasi file \tekskode{TemplateTugas.tex} pada folder \tekskode{Non IUP}
     \item Jawaban atas pertanyaan tugas \textbf{SEMUA} diletakkan pada file \tekskode{assignmentanswers.tex} (untuk IUP) atau \tekskode{jawaban.tex} (untuk non IUP)
    \end{enumerate}
     \item Pekerjaan dikumpulkan dalam bentuk file PDF hasil kompilasi kode LaTeX yang diberikan. Kode LaTeX (project dala format ZIP) juga dilampirkan.
     \item Tugas dikumpulkan selambat-lambatnya tanggal 
     \begin{tcolorbox}
     \center
     \tekskode{12 Oktober 2026 jam 11.00 WIB}    
     \end{tcolorbox}
     \textbf{Pengumpulan yang terlambat akan ditolak oleh sistem dan tidak difasilitasi pengumpulan tugas menggunakan platform lainnya (email, chat, WA, dsb.)}
     \item SETIAP anggota kelompok perlu submit tugas secara individu (tidak hanya perwakilan kelompok).
     \item Setiap pengumpulan tugas HARUS memuat identitas anggota kelompok secara lengkap dan jelas
     \item Pengumpulan melalui Google Classroom harus disertai dengan menekan tombol \textbf{Submit} dan mendapatkan informasi pengiriman/feedback dari sistem.
\end{enumerate}
\end{warnbox}
\end{document}
